\section{La sucesión de Mayer--Vietoris}

% ------------------------------------------------------------------------
\subsection{Enunciado}

\begin{frame}
  \frametitle{La sucesión de Mayer--Vietoris}

  Sean $A,B\subseteq X$ con $\Int A\cup\Int B=X$, y sean
  \[
    i_A\colon A\cap B\hookrightarrow A,\quad
    i_B\colon A\cap B\hookrightarrow B,\quad
    j_A\colon A\hookrightarrow X,\quad
    j_B\colon B\hookrightarrow X
  \]
  las inclusiones.

  \begin{theorem}[Mayer--Vietoris]
    Existe una sucesión exacta larga
    \begin{multline*}
      \cdots\to H_n(A\cap B)\xrightarrow{\left(\begin{smallmatrix}{i_A}_*\\-{i_B}_*\end{smallmatrix}\right)}H_n(A)\oplus H_n(B)
      \xrightarrow{\left(\begin{smallmatrix}{j_A}_* & {j_B}_*\end{smallmatrix}\right)}H_n(X)\\
      \xrightarrow{\ \delta\ }H_{n-1}(A\cap B)\to\cdots\to H_0(X)\to0.
    \end{multline*}
  \end{theorem}

  Este teorema es consecuencia del lema siguiente y del teorema de las
  cadenas pequeñas.
\end{frame}

% ------------------------------------------------------------------------
\subsection{Demostración}

\begin{frame}[allowframebreaks]
  \frametitle{Una sucesión exacta corta de complejos}

  Sea $\cU=\{A,B\}$, de modo que
  $\CU_\bullet(X)=\Sing_\bullet(A)+\Sing_\bullet(B)$.

  \begin{lemma}
    La sucesión de complejos
    \[
      0\to\Sing_\bullet(A\cap B)\xrightarrow{\left(\begin{smallmatrix}{i_A}_\#\\-{i_B}_\#\end{smallmatrix}\right)}
      \Sing_\bullet(A)\oplus\Sing_\bullet(B)\xrightarrow{\left(\begin{smallmatrix}{j_A}_\# & {j_B}_\#\end{smallmatrix}\right)}
      \CU_\bullet(X)\to0
    \]
    es exacta.
  \end{lemma}

  \begin{proof}
    Obvio porque $\CU_\bullet(X)=\Sing_\bullet(A)+\Sing_\bullet(B)$ y $\Sing_\bullet(A)\cap\Sing_\bullet(B)=\Sing_\bullet(A\cap B)$.
  \end{proof}
\end{frame}

\begin{frame}
  \frametitle{El morfismo de conexión}

  \begin{proposition}
    Sea $\alpha\in H_n(X)$. Existe un ciclo $x+y$ que lo representa con
    $x\in\Sing_n(A)$ e $y\in\Sing_n(B)$, y para cualquiera de ellos
    \[
      \delta\alpha=[\partial x]\in H_{n-1}(A\cap B).
    \]
  \end{proposition}

  \begin{proof}
    Por el teorema de las cadenas pequeñas, $\alpha$ tiene un representante
    en $\CU_n(X)=\Sing_n(A)+\Sing_n(B)$. Por la construcción del morfismo de
    conexión, se eleva $x+y\in \CU_n(X)$ a $(x,y)\in \Sing_n(A)\oplus\Sing_n(B)$, se toma su borde $(\partial
    x,\partial y)$ y se busca su preimagen por $\left(\begin{smallmatrix}{i_A}_\#\\-{i_B}_\#\end{smallmatrix}\right)$. Como
    $\partial x+\partial y=\partial(x+y)=0$, la cadena
    $\partial x=-\partial y$ está en $\Sing_{n-1}(A\cap B)$ y su imagen es
    $(\partial x,-\partial x)=(\partial x,\partial y)$.
  \end{proof}
\end{frame}

% ------------------------------------------------------------------------
\subsection{Versión reducida}

\begin{frame}[allowframebreaks]
  \frametitle{Mayer--Vietoris en homología reducida}

  \begin{theorem}
    Si además $A\cap B\neq\varnothing$, hay una sucesión exacta larga
    \begin{multline*}
      \cdots\to\widetilde{H}_n(A\cap B)\xrightarrow{\left(\begin{smallmatrix}{i_A}_*\\-{i_B}_*\end{smallmatrix}\right)}
      \widetilde{H}_n(A)\oplus\widetilde{H}_n(B)
      \xrightarrow{\left(\begin{smallmatrix}{j_A}_* & {j_B}_*\end{smallmatrix}\right)}\widetilde{H}_n(X)\\
      \xrightarrow{\ \delta\ }\widetilde{H}_{n-1}(A\cap B)\to\cdots
      \to\widetilde{H}_0(X)\to0,
    \end{multline*}
    con el morfismo de conexión $\delta$ dado por la misma fórmula.
  \end{theorem}

  \framebreak

  \begin{proof}
    Sea $\tCU_\bullet(X)$ el complejo aumentado
    de las cadenas subordinadas. La sucesión del lema se extiende a los
    complejos aumentados:
    \[
      0\to\widetilde{\Sing}_\bullet(A\cap B)\xrightarrow{\left(\begin{smallmatrix}{i_A}_\#\\-{i_B}_\#\end{smallmatrix}\right)}
      \widetilde{\Sing}_\bullet(A)\oplus\widetilde{\Sing}_\bullet(B)
      \xrightarrow{\left(\begin{smallmatrix}{j_A}_\# & {j_B}_\#\end{smallmatrix}\right)}\tCU_\bullet(X)\to0.
    \]
    En grado $-1$ es $0\to\Z\to\Z\oplus\Z\to\Z\to0$, con $x\mapsto(x,-x)$ y
    $(y,z)\mapsto y+z$, que es exacta.

    \medskip

    Su sucesión exacta larga, junto con la versión reducida del teorema de
    las cadenas pequeñas, da la sucesión del enunciado. Termina en
    $\widetilde{H}_0(X)\to0$ porque
    $\widetilde{H}_{-1}(A\cap B)=0$, ya que $A\cap B\neq\varnothing$.
  \end{proof}
\end{frame}
